% !TEX program = xelatex
\IfFileExists{main.tex}{\documentclass[main.tex]{subfiles}}{
  \IfFileExists{../../main.tex}{\documentclass[../../main.tex]{subfiles}}{
    \IfFileExists{../main.tex}{\documentclass[../main.tex]{subfiles}}{
      \documentclass[Book\space NL/main.tex]{subfiles}
    }
  }
}
\begin{document}
\ifSubfilesClassLoaded{
  \ifcsname appendixchapterparent\endcsname
  \else
    \standaloneappendixchaptersetup{C}{4}
  \fi
}{}
\chapter{Introduction to Medical Parasitology}
\margintoc{}

\textbf{Parasitism}\index{parasitism} is defined as an interaction between an organism which depends on a host for nutrients, and the host, who is harmed by the organism. Consequently \textbf{parasitology}\index{parasitology} is the study of parasites and their interactions with the host. Medical parasitology focuses on the medical implications of parasites that infect humans. This includes: morphology, life-cycle, epidemiology, pathology, diagnosis, prevention and control of parasites.\footnote{Obviously not mind control of course. I hope that this is abundantly evident to the reader.} Parasites are classified into two major groups by their characteristics: unicellular parasites and multicellular parasites. Multicellular parasites are further classified into helminths (nematodes, trematodes and cestodes) and arthropods (fleas, ticks, lice, mosquitoes). Generally, the infective site of parasites may be divided into either gastrointestinal or blood an tissue. Notably, some parasites may occur \textbf{ectopically} (meaning: infecting outside of their typical site) in many cases.

The lifecycle of a parasite begins with infection of the host by the \textbf{infective stage}\index{infective stage (of parasite)}, which may be through a variety of mechanisms, including the ingestion of undercooked food, vector-bourne transmission, skin penetration, sexual transmission and the fecal-oral route. Following infection, the parasite establishes an infection and develops. The parasite then releases progeny (which may also be the \textbf{diagnostic stage}\index{diagnostic stage (of parasite)} that is detectable in fomites produced by the host) that infects organisms or other hosts. At any point, the stage which causes disease in the host is referred to as the \textbf{pathogenic stage}\index{pathogenic stage (of parasite)}.

Laboratory diagnosis of parasites typically involves microscopy (definitive identification), immunological and molecular techniques, though others may be used in unique parasites. Specimens depend on a pathogen's infective site, though common specimens are stool, blood (plasma, RBC, WBC), sputum/bronchoalveolar lavage, CSF, urine, urogenital secretions (vaginal discharge, pus) and tissue (biopsy/autopsy).

Key term definitions follow:

location:
\begin{itemize}
  \item Ectoparasite\index{ectoparasite}\index{parasite!ectoparasite} --- a parasite that lives on the body surface or in the superficial tissues of its host (e.g.\ lice, ticks, mites). Its presence is called an \textit{infestation}\index{infestation}.
  \item Endoparasite\index{endoparasite}\index{parasite!endoparasite} --- a parasite that lives inside its host, in tissues, organs, body cavities, blood or cells (e.g.\ \textit{Plasmodium}, \textit{Ascaris}). Its presence is called an \textit{infection}\index{infection}.
\end{itemize}

Ability to live outside the host:
\begin{itemize}
  \item Facultative parasite\index{facultative parasite} --- an organism that can live either free-living or as a parasite, depending on conditions.
  \item Obligatory parasite\index{obligatory parasite}\index{parasite!obligatory} --- a parasite that cannot survive or complete its life cycle without a host (e.g.\ \textit{Plasmodium}, \textit{Toxoplasma gondii}).
  \item Accidental parasite\index{accidental parasite} --- a parasite found in a host it does not normally infect.
\end{itemize}

Types of host:
\begin{itemize}
  \item Definitive host\index{definitive host}\index{host!definitive} --- the host where the parasite reaches its adult stage and potentially reproduces sexually.
  \item Intermediate host\index{intermediate host} --- the host where the parasite's larval or asexual stages develop. Some parasites need more than one, called the first and second intermediate hosts.
  \item Reservoir host\index{reservoir host}\index{host!reservoir} --- an animal that carries a parasite that also infects humans and acts as a continuing source of human infection.
  \item Accidental host\index{accidental host}\index{host!accidental}\index{dead-end host|see{accidental host}} --- a host that is not part of the parasite's normal life cycle. It is usually a dead end, so the parasite cannot develop further or pass to another host.
  \item Paratenic host\index{paratenic host} (transport host) --- host in which the immature parasite does not develop.
\end{itemize}

Periods of infection
\begin{itemize}
  \item Pre-patent period\index{pre-patent period}\index{infection!pre-patent period} --- the time from infection until the parasite's diagnostic stages (eggs, larvae, cysts or the parasite itself) can first be detected in blood, faeces or other specimens. It is also called the biological incubation period.
  \item Incubation period\index{incubation period} --- the time from infection until the first clinical signs or symptoms appear. It is also called the clinical incubation period.
  \item Patent period\index{patent period} --- the period when the parasite's diagnostic stages can be detected in the host.
  \item Parthenogenesis\footnote{It is really spelled like this! I did not make a mistake here.} --- parasite asexual reproduction (typically outside of the definitive host).
\end{itemize}

\end{document}